% arXiv-ready LaTeX with bibliography
% External-FDT: Cloud-to-Building Control System Tunneling in BACnet Networks
% cs.CR - Cryptography and Security

% arXiv metadata
%arXiv:cs.CR/2505.00000 [cs.CR]
%Comments: 3 pages, 2 tables. Empirical security analysis of internet-exposed building automation systems.
%Primary Class: cs.CR
%Secondary Class: cs.SY (Systems and Control)

\documentclass[11pt,a4paper]{article}

% Core packages
\usepackage[utf8]{inputenc}
\usepackage[T1]{fontenc}
\usepackage{lmodern}
\usepackage{microtype}
\usepackage{amsmath,amssymb,amsfonts}
\usepackage[numbers,sort&compress]{natbib}
\bibliographystyle{unsrtnat}
\usepackage{hyperref}
\usepackage{url}
\usepackage{booktabs}
\usepackage{graphicx}
\usepackage[margin=1in]{geometry}
\usepackage{xcolor}

\hypersetup{
  pdfauthor={Cameron Warren},
  pdftitle={External-FDT: Cloud-to-Building Control System Tunneling in BACnet Networks},
  colorlinks=true,
  linkcolor=blue,
  citecolor=blue
}

% Title
\title{External-FDT: Cloud-to-Building Control System Tunneling in BACnet Networks}
\author{Cameron Warren\\
  \texttt{cwarre33@uncc.edu}\\[0.5em]
  Department of Computer Science\\
  University of North Carolina at Charlotte
}
\date{\today}

\begin{document}

\maketitle

\begin{abstract}
Building automation systems increasingly bridge private facility networks to public 
cloud infrastructure, creating undocumented attack surfaces. Through longitudinal 
analysis of internet-exposed BACnet Broadcast Management Devices (BBMDs), we 
discover a systemic pattern we term \textbf{External-FDT}: 71\% of publicly visible 
BBMDs maintain Foreign Device Table entries pointing to public, routable IP addresses 
rather than internal network hosts. Of these, 42\% tunnel to major cloud providers 
(AWS, Azure, DigitalOcean), while 25\% reveal a novel \textbf{integrator concentration} 
phenomenon where a single external endpoint serves multiple geographically separated 
building networks. Our empirical analysis of 17 seeded BBMDs (longitudinally sampled over 
47 days) characterizes this behavior and exposes potential supply-chain attack 
surfaces in building automation infrastructure. All findings derive from passive 
OSINT analysis of Shodan's indexed banners---no active scanning of target systems.
\end{abstract}

\section{Introduction}
\label{sec:intro}

Building automation systems---responsible for HVAC, access control, fire safety, and 
critical facility operations---have historically relied on isolated 
broadcast networks. The BACnet protocol (ISO 16484-5)\cite{ashrae2016bacnet} enables these systems 
to span IP subnets through Broadcast Management Devices (BBMDs) that forward 
broadcast packets to registered Foreign Devices (FDs). These tunnels are 
typically confined to private networks\cite{newman2009bacnet}.

We discover a systemic deviation from this model.

\subsection{The External-FDT Phenomenon}

Our analysis of internet-exposed BBMDs reveals that the majority maintain 
Foreign Device Table (FDT) entries pointing to \textbf{public, routable IP addresses}---
creating undocumented bridges between building control networks and external 
infrastructure including major cloud providers (AWS, Azure, DigitalOcean).

This finding is distinct from previously characterized ICS exposure patterns.
Prior work has focused on open ports\cite{tank2017sok}, protocol vulnerabilities\cite{bradley2020industrial}, 
and device enumeration\cite{shodanics}, but no study has characterized \emph{tunneling relationships}
between building control systems and external infrastructure.

\subsection{Key Contributions}

This paper makes three empirical contributions:

\begin{enumerate}
  \item \textbf{Characterization of External-FDT behavior}: We quantify the 
    prevalence (71\%) and patterns (cloud vs. integrator vs. unknown) of 
    public-IP Foreign Device registrations in exposed BBMDs.
    
  \item \textbf{Discovery of Integrator Concentration}: We identify a 
    novel pattern where a single external endpoint serves multiple 
    distinct building networks, suggesting supply-chain attack surfaces.
    
  \item \textbf{Attribution Framework}: We validate passive OSINT methods for 
    attributing external endpoints to cloud providers through IP-to-ASN 
    mapping without active probing.
\end{enumerate}

\section{Background}
\label{sec:background}

\subsection{BACnet and BBMDs}

BACnet (Building Automation and Control Networks) is a standard protocol 
for intelligent building systems. BBMDs (Broadcast Management Devices) 
enable BACnet to span IP subnets by:\cite{newman2009bacnet}

\begin{enumerate}
  \item Receiving broadcast messages on their local subnet
  \item Forwarding encapsulated broadcasts to registered Foreign Devices
  \item Receiving unicasts from Foreign Devices and re-broadcasting locally
\end{enumerate}

A Foreign Device Table entry contains: IP address, port, TTL, and registration 
timestamp. Critically, the FD IP need not be on the same network as the BBMD---
making cloud-to-building bridging architecturally possible.

\subsection{Related Work}

Prior ICS security research has focused on:
\begin{itemize}
  \item Protocol vulnerabilities and specific CVE exploitation\cite{bradley2020industrial}
  \item Shodan-based exposure quantification\cite{tank2017sok,shodanics}
  \item Modbus, DNP3, and BACnet object enumeration\cite{gamage2018security}
\end{itemize}

To our knowledge, no prior work has characterized \emph{tunneling relationships}
between internet-exposed building control systems and external infrastructure.

\section{Methodology}
\label{sec:methodology}

\subsection{Ethical Constraints}

All analysis uses \textbf{passive OSINT only}: Shodan-indexed banners captured 
during normal internet scanning\cite{shodanics}. We performed:
\begin{itemize}
  \item No active connection attempts to target IPs
  \item No BACnet Who-Is or Read-Property requests
  \item No authentication attempts
  \item No packet transmission to target networks
\end{itemize}

\subsection{Data Source}

Longitudinal Shodan dataset: 17 seeded BBMDs known to exhibit BACnet/WP 
fingerprinting, sampled March 4, 2026 through April 20, 2026 (47 days), 
enriched with Foreign Device Table extractions.

\subsection{Classification Schema}

We classify each Foreign Device Table entry:

\begin{table}[h]
\centering
\begin{tabular}{lll}
\toprule
\textbf{Category} \& \textbf{FD IP Range} \& \textbf{Interpretation} \\
\midrule
Internal-FDT \& RFC1918 (10/8, 172.16/12, 192.168/16) \& Normal operation \\
External-FDT \& Public routable \& Requires investigation \\
\quad Cloud \& AWS, Azure, DO, GCP IP ranges \& Cloud-hosted management \\
\quad Integrator \& Same FD serves multiple BBMDs \& Third-party monitoring \\
\quad Unknown \& Uncategorized external \& Further analysis needed \\
\bottomrule
\end{tabular}
\caption{FDT Entry Classification}
\label{tab:classification}
\end{table}

Cloud attribution uses IP-to-ASN mapping against known cloud provider 
Autonomous System Numbers\cite{tank2017sok}.

\section{Results}
\label{sec:results}

\subsection{Prevalence of External-FDT Behavior}

Of 17 analyzed BBMDs:
\begin{itemize}
  \item \textbf{12 (71\%)} exhibited External-FDT behavior (public FD IPs)
  \item \textbf{5 (29\%)} maintained only Internal-FDT (RFC1918) entries
\end{itemize}

This ratio suggests External-FDT is \emph{systemic} rather than anomalous.

\subsection{Cloud Provider Distribution}

Among External-FDT cases ($N=12$), destination infrastructure shows 
cloud concentration:

\begin{table}[h]
\centering
\begin{tabular}{lcc}
\toprule
\textbf{Provider} \& \textbf{Count} \& \textbf{Percentage} \\
\midrule
DigitalOcean \& 4 \& 33\% \\
Amazon AWS \& 2 \& 17\% \\
Microsoft Azure \& 1 \& 8\% \\
Other / Unknown \& 5 \& 42\% \\
\midrule
\textbf{Total Cloud} \& 5 \& 42\% \\
\bottomrule
\end{tabular}
\caption{Cloud Provider Attribution (External-FDT Cases)}
\label{tab:clouds}
\end{table}

DigitalOcean shows unexpected dominance---suggesting either 
preferred vendor relationships or lower barriers to entry for 
building automation integrators.

\subsection{Case Study: Integrator Concentration}

The most significant finding: two distinct BBMDs share a common 
Foreign Device IP.

\begin{table}[h]
\centering
\begin{tabular}{ll}
\toprule
\textbf{Attribute} \& \textbf{Value} \\
\midrule
BBMD 1 IP \& 66.58.248.125 \\
BBMD 2 IP \& 24.237.132.230 \\
Shared FD IP \& 216.67.73.166 \\
Geographic separation \& ~4,000 km (different cities) \\
Observation window \& 47 days \\
Port behavior \& 13 rotating source ports each \\
\bottomrule
\end{tabular}
\caption{Integrator Concentration Pattern}
\label{tab:concentration}
\end{table}

\textbf{Implication}: If the owner of 216.67.73.166 is a building 
automation integrator, this single endpoint bridges to at least 
two distinct client facility networks---suggesting a supply-chain 
attack surface where compromising one monitoring station yields 
access to multiple building control systems.

\subsection{Persistence Analysis}

Longitudinal data shows External-FDT relationships are stable:
\begin{itemize}
  \item Average observation count: 47 scans per relationship
  \item Maximum: 930 observations over 47 days (single case)
  \item Port rotation patterns suggest NAT vs. direct connection
\end{itemize}

\section{Discussion}
\label{sec:discussion}

\subsection{Why This Matters}

External-FDT creates three attack surfaces:

\begin{enumerate}
  \item \textbf{Cloud compromise}: Breach of cloud-hosted management 
    infrastructure yields direct access to building control networks
  \item \textbf{Integrator pivot}: Single integrator compromise affects 
    multiple facilities
  \item \textbf{Exposure amplification}: Internet exposure of the FD 
    endpoint compounds building network exposure
\end{enumerate}

\subsection{Limitations}

Passive analysis cannot distinguish:
\begin{itemize}
  \item Authorized SaaS integration vs. unauthorized bridging
  \item Active exploitation vs. dormant configuration
  \item Bidirectional tunnels vs. unidirectional monitoring
\end{itemize}

Attribution of cloud IPs to specific tenants or operators requires 
additional investigation.

\section{Conclusion}

We empirically demonstrate that External-FDT---cloud-to-building BACnet 
tunneling---is systemic in internet-exposed building automation infrastructure. 
The discovery of integrator concentration patterns suggests novel supply-chain 
attack surfaces requiring further research and coordinated disclosure.

\section*{Acknowledgments}

This research was conducted using only publicly available Shodan index data.
All facility identities and IP addresses have been anonymized except where 
already public in prior disclosures.

\bibliography{references}

\end{document}
